\documentclass[11pt]{article}
\usepackage{amsmath, amssymb, amsthm}
\usepackage{mathpartir}
\usepackage{listings}
\usepackage{xcolor}

\newcommand{\syn}{\Rightarrow} % Synthesis/inference
\newcommand{\chk}{\Leftarrow} % Checking

\title{Dahlia Type System}
\author{Charlotte W.}
\date{\today}

\begin{document}
\maketitle

\section{Preamble}
This document defines the formal type system for the Dahlia language.
Dahlia is a tiny low-level, memory-safe programming language utilizing a static type system with type inference based on Hindley-Milner and Algorithm W.

\section{Data Types}
\begin{itemize}
    \item \textbf{i8-i64}: A sized integer type (\texttt{i8-i64}) with specified bit width. Unsigned integers are denoted with a \texttt{u} rather than an \texttt{i} prefix.
    \item \textbf{bool}: A boolean type (\texttt{bool}) representing true or false values.
    \item \textbf{char}: A single character type (\texttt{char}) (8 bits wide).
    \item \textbf{str}: A string type (\texttt{str}) representing a sequence of characters.
    \item \textbf{Str}: A boxed string type (\texttt{Str}) representing a heap-allocated string.
    \item \textbf{Array}: A fixed-size array type (\texttt{[T; N]}) representing a collection of elements of type \texttt{T} with size \texttt{N}.
\end{itemize}


\section{Constants}
Constants are used to define values that cannot be changed during program execution.
They are defined using: \texttt{const <name>: <type> = <value>}

\section{How to Read the Rules}
Each rule has one or more premises above the line and a conclusion below it. If every premise holds, then the conclusion is valid.

\begin{itemize}
    \item $\Gamma$ is the typing context, mapping names to types.
    \item $\tau$ and $\tau_r$ range over types.
    \item $e$ ranges over expressions, and $b$ ranges over blocks of statements.
    \item $\syn$ means a form synthesizes its type from the term itself.
    \item $\chk$ means a form is checked against an expected type.
    \item $\Rightarrow$ is used for statements and declarations that update or preserve the context.
    \item Items listed as "for each" or written with subscripts such as $e_i$ must hold for every matching element.
\end{itemize}

\section{Generic Type Syntax}
Polymorphic types are written with square-bracket type parameters, for example \texttt{List[A]} or \texttt{Result[A, E]}. Type constructors are written in the constructor position, followed by their type arguments in square brackets. A type scheme is written as $\forall \alpha_1, \ldots, \alpha_n.\,\tau$, where the quantified variables may be instantiated with fresh types during lookup.

\begin{mathpar}
    \inferrule* [lab=TyApp]
    {\vdash \tau_1 \text{ type} \\
        \cdots \\
        \vdash \tau_n \text{ type} \\
        C \text{ is a type constructor of arity } n}
    {\vdash C[\tau_1, \ldots, \tau_n] \text{ type}}

    \and

    \inferrule* [lab=TyScheme]
    {\vdash \tau \text{ type} \\
        \alpha_1, \ldots, \alpha_n \notin \mathrm{ftv}(\Gamma)}
    {\Gamma \vdash \forall \alpha_1, \ldots, \alpha_n.\,\tau \text{ scheme}}
\end{mathpar}

This means that \texttt{List[A]} is well formed when \texttt{List} expects one type argument and \texttt{A} is itself a valid type, and similarly for constructors with more parameters such as \texttt{Result[A, E]}. A type scheme such as $\forall A.\,List[A] \to A$ is not a distinct runtime type; it is a polymorphic binding that may later be instantiated with fresh type variables.

\section{Type Checking}
The type system employs bidirectional type checking with synthesis and checking.
The type system itself is based on Hindley-Milner, allowing for polymorphism and type inference.

Blocks are typed by the value of their terminal expression or explicit return; an empty block has type \texttt{void}.

\begin{itemize}
    \item \textbf{Synthesis}: $\Gamma \vdash e \syn \tau$ means an expression synthesizes its type from the term in $\Gamma$.
    \item \textbf{Checking}: $\Gamma \vdash e \chk \tau$ means an expression is checked against an expected type $\tau$ in $\Gamma$.
\end{itemize}

The rules below are aligned with \texttt{GRAMMAR.md} and the current AST. Dahlia does not currently have lambda syntax or explicit type annotations on expressions, so the previous $\textsc{Ann}$ and $\textsc{Abs}$ rules do not apply.

\paragraph{Literal and variable rules}
\begin{mathpar}
    \inferrule* [lab=Lit]
    { }
    {\Gamma \vdash \ell \syn \tau_{\ell}}

    \and

    \inferrule* [lab=Var]
    {x : \tau \in \Gamma}
    {\Gamma \vdash x \syn \tau}

    \and

    \inferrule* [lab=ArrayLit]
    {\Gamma \vdash e_i \chk \tau \text{ for each } i \in 1..n}
    {\Gamma \vdash [e_1, \ldots, e_n] \syn [\tau; n]}
\end{mathpar}

\paragraph{Expression rules}
\begin{mathpar}
    \inferrule* [lab=Call]
    {\Gamma \vdash f \syn (\tau_1, \ldots, \tau_n) \to \tau_r \\
        \Gamma \vdash e_i \chk \tau_i \text{ for each } i \in 1..n}
    {\Gamma \vdash f(e_1, \ldots, e_n) \syn \tau_r}

    \and

    \inferrule* [lab=If]
    {\Gamma \vdash c \chk \texttt{bool} \\
        \Gamma \vdash b_1 \syn \tau \\
        \Gamma \vdash b_2 \syn \tau}
    {\Gamma \vdash \texttt{if } c \{ b_1 \} \texttt{ else } \{ b_2 \} \syn \tau}

    \and

    \inferrule* [lab=Addr]
    {\Gamma \vdash e \chk \tau}
    {\Gamma \vdash \&e \syn *\tau}

    \and

    \inferrule* [lab=Deref]
    {\Gamma \vdash e \syn *\tau}
    {\Gamma \vdash *e \syn \tau}
\end{mathpar}

\paragraph{Operator rules}
\begin{mathpar}
    \inferrule* [lab=Arith]
    {\Gamma \vdash e_1 \chk \tau \\
        \Gamma \vdash e_2 \chk \tau \\
        \tau \in \{\texttt{i8}, \texttt{i16}, \texttt{i32}, \texttt{i64}, \texttt{u8}, \texttt{u16}, \texttt{u32}, \texttt{u64}, \texttt{f32}, \texttt{f64}\}}
    {\Gamma \vdash e_1 \; \oplus \; e_2 \syn \tau}

    \and

    \inferrule* [lab=Cmp]
    {\Gamma \vdash e_1 \chk \tau \\
        \Gamma \vdash e_2 \chk \tau \\
        \tau \in \{\texttt{i8}, \texttt{i16}, \texttt{i32}, \texttt{i64}, \texttt{u8}, \texttt{u16}, \texttt{u32}, \texttt{u64}, \texttt{f32}, \texttt{f64}\}}
    {\Gamma \vdash e_1 \; \diamond \; e_2 \syn \texttt{bool}}

    \and

    \inferrule* [lab=Eq]
    {\Gamma \vdash e_1 \chk \tau \\
        \Gamma \vdash e_2 \chk \tau}
    {\Gamma \vdash e_1 \; == \; e_2 \syn \texttt{bool}}

    \and

    \inferrule* [lab=Unary]
    {\Gamma \vdash e \chk \tau}
    {\Gamma \vdash -e \syn \tau}

    \and

    \inferrule* [lab=Not]
    {\Gamma \vdash e \chk \texttt{bool}}
    {\Gamma \vdash !e \syn \texttt{bool}}
\end{mathpar}

\paragraph{Statement rules}
\begin{mathpar}
    \inferrule* [lab=VarDecl]
    {\Gamma \vdash e \chk \tau}
    {\Gamma \vdash \texttt{var } x : \tau = e \Rightarrow \Gamma, x : \tau}

    \and

    \inferrule* [lab=ConstDecl]
    {\Gamma \vdash e \chk \tau}
    {\Gamma \vdash \texttt{const } x : \tau = e \Rightarrow \Gamma, x : \tau}

    \and

    \inferrule* [lab=Assign]
    {x : \tau \in \Gamma \\
        \Gamma \vdash e \chk \tau}
    {\Gamma \vdash x = e \Rightarrow \Gamma}

    \and

    \inferrule* [lab=While]
    {\Gamma \vdash c \chk \texttt{bool} \\
        \Gamma \vdash b \Rightarrow \Gamma}
    {\Gamma \vdash \texttt{while } c \{ b \} \Rightarrow \Gamma}

    \and

    \inferrule* [lab=For]
    {\Gamma \vdash it \syn [\tau; n] \\
        \Gamma, x : \tau \vdash b \Rightarrow \Gamma}
    {\Gamma \vdash \texttt{for } x : it \{ b \} \Rightarrow \Gamma}

    \and

    \inferrule* [lab=Return]
    {\Gamma \vdash e \chk \tau_r}
    {\Gamma \vdash \texttt{return } e \Rightarrow \Gamma}

    \and

    \inferrule* [lab=Break]
    { }
    {\Gamma \vdash \texttt{break} \Rightarrow \Gamma}
\end{mathpar}

\paragraph{Declarations and polymorphism}
\begin{mathpar}
    \inferrule* [lab=Inst]
    {x : \forall \alpha_1, \ldots, \alpha_n.\,\tau \in \Gamma}
    {\Gamma \vdash x \syn \tau[\alpha_1 \mapsto \beta_1, \ldots, \alpha_n \mapsto \beta_n]}

    \and

    \inferrule* [lab=Gen]
    {\Gamma \vdash e \syn \tau \\
        \vec{\alpha} = \mathrm{ftv}(\tau) \setminus \mathrm{ftv}(\Gamma)}
    {\Gamma \vdash e \rightsquigarrow \forall \vec{\alpha}.\,\tau}

    \and

    \inferrule* [lab=VarDecl]
    {\Gamma \vdash e \chk \tau}
    {\Gamma \vdash \texttt{var } x : \tau = e \Rightarrow \Gamma, x : \tau}

    \and

    \inferrule* [lab=ConstDecl]
    {\Gamma \vdash e \syn \tau \\
        \sigma = \mathrm{Gen}(\Gamma, \tau)}
    {\Gamma \vdash \texttt{const } x : \tau = e \Rightarrow \Gamma, x : \sigma}

    \and

    \inferrule* [lab=FnDecl]
    {\Gamma, x_1 : \tau_1, \ldots, x_n : \tau_n \vdash b \syn \tau_r \\
        \sigma = \mathrm{Gen}(\Gamma, \tau_1 \to \cdots \to \tau_n \to \tau_r)}
    {\Gamma \vdash \texttt{fn } f(x_1 : \tau_1, \ldots, x_n : \tau_n) \{ b \} \Rightarrow \Gamma, f : \sigma}
\end{mathpar}

\begin{itemize}
    \item \textbf{FnDecl}: function parameters are monomorphic within the body; the function name may be generalized after the body type is known.
    \item \textbf{VarDecl}: mutable bindings remain monomorphic so later assignments preserve a single type.
    \item \textbf{ConstDecl}: immutable bindings may be generalized because their type does not change after binding.
    \item \textbf{StructDecl}: struct fields must have unique names and well-formed types.
    \item \textbf{EnumDecl}: enum variants must have unique names.
    \item \textbf{Generic types}: type constructors use square-bracket syntax such as \texttt{List[A]}; the constructor arity must match the number of supplied arguments.
    \item \textbf{Instantiate}: each lookup of a polymorphic binding replaces quantified variables with fresh type variables.
    \item \textbf{Generalize}: when binding a value, quantify only type variables that are free in the inferred type but not free in the environment.
    \item \textbf{Unify}: synthesis and checking meet through unification with occurs-check.
\end{itemize}

\end{document}
